\section{Termodinamica}
\subsection{Trasformazioni}

\begin{longtable}[]{@{}lccccc@{}}
\toprule
\begin{minipage}[b]{0.07\columnwidth}\raggedright
Tipo\strut
\end{minipage} & \begin{minipage}[b]{0.13\columnwidth}\centering
\(W\)\strut
\end{minipage} & \begin{minipage}[b]{0.09\columnwidth}\centering
\(Q\)\strut
\end{minipage} & \begin{minipage}[b]{0.09\columnwidth}\centering
\(\Delta U\)\strut
\end{minipage} & \begin{minipage}[b]{0.39\columnwidth}\centering
Costanti\strut
\end{minipage} & \begin{minipage}[b]{0.05\columnwidth}\centering
\(n\)\strut
\end{minipage}\tabularnewline
\midrule
\endhead
\begin{minipage}[t]{0.07\columnwidth}\raggedright
Politropica\strut
\end{minipage} & \begin{minipage}[t]{0.13\columnwidth}\centering
\strut
\end{minipage} & \begin{minipage}[t]{0.09\columnwidth}\centering
\strut
\end{minipage} & \begin{minipage}[t]{0.09\columnwidth}\centering
\strut
\end{minipage} & \begin{minipage}[t]{0.39\columnwidth}\centering
\(pV^n\)\strut
\end{minipage} & \begin{minipage}[t]{0.05\columnwidth}\centering
\(n\)\strut
\end{minipage}\tabularnewline
\begin{minipage}[t]{0.07\columnwidth}\raggedright
Isocora\strut
\end{minipage} & \begin{minipage}[t]{0.13\columnwidth}\centering
0\strut
\end{minipage} & \begin{minipage}[t]{0.09\columnwidth}\centering
\(nc_v\Delta T\)\strut
\end{minipage} & \begin{minipage}[t]{0.09\columnwidth}\centering
\(Q\)\strut
\end{minipage} & \begin{minipage}[t]{0.39\columnwidth}\centering
\(\frac{p}{T}\)\strut
\end{minipage} & \begin{minipage}[t]{0.05\columnwidth}\centering
\(\infty\)\strut
\end{minipage}\tabularnewline
\begin{minipage}[t]{0.07\columnwidth}\raggedright
Isobara\strut
\end{minipage} & \begin{minipage}[t]{0.13\columnwidth}\centering
\(p\Delta V\)\strut
\end{minipage} & \begin{minipage}[t]{0.09\columnwidth}\centering
\(nc_p\Delta T\)\strut
\end{minipage} & \begin{minipage}[t]{0.09\columnwidth}\centering
\(nc_v\Delta T\)\strut
\end{minipage} & \begin{minipage}[t]{0.39\columnwidth}\centering
\(\frac{V}{T}\)\strut
\end{minipage} & \begin{minipage}[t]{0.05\columnwidth}\centering
\(0\)\strut
\end{minipage}\tabularnewline
\begin{minipage}[t]{0.07\columnwidth}\raggedright
Isoterma\strut
\end{minipage} & \begin{minipage}[t]{0.13\columnwidth}\centering
\(nRT\ln{(V_f/V_i)}\)\strut
\end{minipage} & \begin{minipage}[t]{0.09\columnwidth}\centering
\(W\)\strut
\end{minipage} & \begin{minipage}[t]{0.09\columnwidth}\centering
0\strut
\end{minipage} & \begin{minipage}[t]{0.39\columnwidth}\centering
\(pV\)\strut
\end{minipage} & \begin{minipage}[t]{0.05\columnwidth}\centering
\(1\)\strut
\end{minipage}\tabularnewline
\begin{minipage}[t]{0.07\columnwidth}\raggedright
Adiabatica\strut
\end{minipage} & \begin{minipage}[t]{0.13\columnwidth}\centering
\(-\Delta U\)\strut
\end{minipage} & \begin{minipage}[t]{0.09\columnwidth}\centering
0\strut
\end{minipage} & \begin{minipage}[t]{0.09\columnwidth}\centering
\(nc_v\Delta T\)\strut
\end{minipage} & \begin{minipage}[t]{0.39\columnwidth}\centering
\(pV^\gamma , Tp^{\frac{1-\gamma}{\gamma}}, TV^{\gamma-1}\)\strut
\end{minipage} & \begin{minipage}[t]{0.05\columnwidth}\centering
\(\gamma\)\strut
\end{minipage}\tabularnewline
\bottomrule
\end{longtable}

Equazione di stato dei gas perfetti: \(pV=nRT\) Energia:
\(U=nc_vT\implies \Delta U = nc_v\Delta T\) Primo principio della
termodinamica: \(\Delta U = Q-W\) Costante dei gas perfetti:
\(R=8.314 \frac{J}{mol\cdot K}\) Costante specifica del gas: \(R^*=R/M\)
(\(\frac{J}{kg\cdot K}\)) Relazione di Meyer: \(R=c_p-c_v\) Costante
gamma: \(\gamma=\frac{c_p}{c_v}\) Calore specifico molare in gas
monoatomici: \(c_v=\frac{3}{2}R\), \(c_p=\frac{5}{2}R\),
\(\gamma=\frac{5}{3}\) Calore specifico molare in gas biatomici:
\(c_v=\frac{5}{2}R\), \(c_p=\frac{7}{2}R\), \(\gamma=\frac{7}{5}\)
Calore specifico molare in gas poliatomici: \(c_v=3R\), \(c_p=4R\),
\(\gamma=\frac{4}{3}\) Lavoro: \[L_{AB}=\int_A^B p(V)dv\] Variazione
entropia:
\[\Delta S_{AB}=nc_p\ln{\left(\frac{T_B}{T_A}\right)}-nR\ln{\left(\frac{p_B}{p_A}\right)}=nc_v\ln{\left(\frac{T_B}{T_A}\right)}+nR\ln{\left(\frac{V_B}{V_A}\right)}\]
Variazione entropia specifica:
\[\Delta s_{AB}=\frac{\Delta S_{AB}}{n}=c_p\ln{\left(\frac{T_B}{T_A}\right)}-R\ln{\left(\frac{p_B}{p_A}\right)}=c_v\ln{\left(\frac{T_B}{T_A}\right)}+R\ln{\left(\frac{V_B}{V_A}\right)}\]
\#\# Macchine e cicli In un ciclo l'energia interna rimane costante:
\(A = B \implies \Delta U= 0\) Macchina termica Lavoro:
\(W = Q_A-|Q_C|\) Rendimento:
\(\eta =\frac{W}{Q_A}=1-\frac{|Q_C|}{Q_A}\) Il rendimento in linea
generale aumenta all'aumentare della differenza di temperatura in
ingresso e in uscita. \#\#\# Ciclo di Carnot Composto da: - Espansione
isoterma (non varia temperatura) - Espansione isoentropica (non varia
entropia) - Compressione isoterma - Compressione isoentropica

Rendimento: \(\eta_\text{Carnot}=1-\frac{T_i}{T_s}\) Coefficiente di
prestazione: \(COP=\frac{T_i}{T_i-T_s}\)
\includegraphics{Ciclo Carnot pv.png}
\includegraphics{Ciclo Carnot ts.png}

\hypertarget{ciclo-otto}{%
\subsubsection{Ciclo Otto}\label{ciclo-otto}}

\begin{enumerate}
\def\labelenumi{\arabic{enumi}.}
\item
  Compressione isoentropica
\item
  Trasformazione isocora
\item
  Espansione isoentropica
\item
  Trasformazione isocora
\end{enumerate}

Rapporto di compressione: \(r=\frac{V_\text{max}}{V_\text{min}}\)
Rendimento: \(\eta_\text{Otto}=1-\frac{1}{r^{k-1}}\)
\includegraphics{Ciclo Otto pv.png} \includegraphics{Ciclo Otto ts.png}

\hypertarget{ciclo-diesel}{%
\subsubsection{Ciclo Diesel}\label{ciclo-diesel}}

\begin{enumerate}
\def\labelenumi{\arabic{enumi}.}
\item
  Compressione isoentropica
\item
  Espansione isobara
\item
  Espansione isoentropica
\item
  Trasformazione isocora \includegraphics{Ciclo Diesel pv.png}
  \includegraphics{Ciclo Diesel ts.png}
\end{enumerate}

\hypertarget{ciclo-joule-brayton}{%
\subsubsection{Ciclo Joule-Brayton}\label{ciclo-joule-brayton}}

\begin{enumerate}
\def\labelenumi{\arabic{enumi}.}
\item
  Compressione isoentropica
\item
  Espansione isobara
\item
  Espansione isoentropica
\item
  Compressione isobara
\end{enumerate}

Rapporto manometrico di compressione: \(\beta=\frac{p_2}{p_1}\)
Coefficiente dilatazione adiabatica: \(k=\frac{c_p}{c_v}\) Rendimento:
\(\eta_\text{Brayton}=1-\frac{1}{\beta^{(k-1)/k}}\) Rendimento massimo
per \(\beta=(T_\text{max}-T_\text{min})^{k/[2(k-1)]}\)
\includegraphics{Ciclo Brayton pv.png}
\includegraphics{Ciclo Brayton ts.png}

\hypertarget{ciclo-di-rankine}{%
\subsubsection{Ciclo di Rankine}\label{ciclo-di-rankine}}

\includegraphics{Ciclo Rankine ts.png}

\hypertarget{macchine-frigorifere-e-pompe-di-calore}{%
\subsubsection{Macchine frigorifere e pompe di
calore}\label{macchine-frigorifere-e-pompe-di-calore}}

\(COP_F=\frac{1}{Q_s/Q_i-1}\hspace{1cm}COP_\text{F,rev}=\frac{1}{T_s/T_i-1}\)
\(COP_\text{PDC}=\frac{1}{1-Q_i/Q_s}\hspace{0.48cm}COP_\text{PDC,rev}=\frac{1}{1-T_i/T_s}\)
\#\#\#\#\# Ciclo inverso a compressione di vapore ideale 1. Compressione
isoentropica 2. Compressione isobara (cede calore) 3. Laminazione
isoentalpica (espansione) 4. Espansione isobara (assorbe calore)

\(COP_F=\frac{q_i}{l_{n,e}}=\frac{h_1-h_4}{h_2-h_1}\)
\(COP_{PDC}=\frac{q_s}{l_{n,e}}=\frac{h_2-h_3}{h_2-h_1}\)
\includegraphics{Ciclo inverso vapore ideale.png}

